
\clearpage
\chapter{דפוסי תיאום מתקדמים}

\section{דפוס הדיבייט (\textenglish{Debate Pattern})}

דפוס \textenglish{Debate} מממש שני סוכנים בעלי נקודות מבט מנוגדות שחייבים להגיע להסכמה לפני שהמחבר מסיים את הטקסט. זה מפחית \textenglish{hallucination} ב-\textenglish{$\sim$31\%} על מטלות ידע (\textenglish{knowledge-intensive tasks}) \cite{segal2024orchestration}.

ניסחנו את הדיבייט כמשחק אפס-סכום חלקי:

\begin{equation}
\text{Consensus}(r_1, r_2) = \begin{cases} r_1 & \text{if } \text{sim}(r_1, r_2) > \theta_c \\ \text{Arbiter}(r_1, r_2) & \text{otherwise} \end{cases}
\end{equation}

כאשר \(\theta_c = 0.82\) (קוסינוס \textenglish{embedding similarity}) וה-\textenglish{Arbiter} הוא מודל \textenglish{LLM} נפרד שמסכם את שתי הדעות.

\section{דפוס \textenglish{Reflection}}

דפוס \textenglish{Reflection} (\textenglish{Self-Refine}) \cite{wu2023autogen} מאפשר לסוכן לבקר את תוצאתו שלו:

\begin{itemize}
\item שלב יצירה: הסוכן מייצר תשובה ראשונית.
\item שלב ביקורת: הסוכן מזהה חולשות בתשובה שלו עצמו.
\item שלב שיפור: יצירה מחדש בהתחשב בביקורת.
\end{itemize}

ניסויים הראו שאחרי \textenglish{2-3} מחזורי \textenglish{reflection}, \textenglish{accuracy} עולה ב-\textenglish{12\%} במשימות קידוד.
